\documentclass[12pt]{article}
\usepackage[utf8]{inputenc}
\usepackage[T1]{fontenc}
\usepackage{hyperref}
\title{Comecrudo Language Resources for K-12 Curriculum Development}
\author{Compiled by Claude Sonnet 5.5 (claude-sonnet-5-5) for Andr\'as Kornai}
\date{}
\begin{document}
\maketitle

\section{Introduction}
Comecrudo (ISO 639-3 xcm), also called Yu\'e, is an extinct language of the lower Rio Grande near Reynosa, Tamaulipas. It is the best-documented member of the small Comecrudan family. Jean-Louis Berlandier recorded a list of 148 words in 1829. Albert Gatschet later recorded vocabulary, sentences and a short text from descendants near Camargo, Tamaulipas, who were no longer fluent. Gatschet's 1886 manuscript is held at the Smithsonian's National Anthropological Archives. Very little else is known. I found no teaching materials, programmes or community language resources.

\section{Primary Educational Resources}
I found no primers, textbooks or classroom materials for Comecrudo.

\section{Linguistic and Reference Materials}
\begin{itemize}
\item Wikipedia, ``Comecrudo language'': overview and sources. \url{https://en.wikipedia.org/wiki/Comecrudo_language}
\item Wikipedia, ``Comecrudan languages'': the family, including Garza and Mamulique. Goddard (1979) argued that they are related. Wider groupings such as Coahuiltecan are disputed. \url{https://en.wikipedia.org/wiki/Comecrudan_languages}
\item Wikipedia, ``Garza language'': a related language known from only 21 words and two tribal names recorded by Berlandier in 1828. \url{https://en.wikipedia.org/wiki/Garza_language}
\end{itemize}
The Berlandier and Gatschet manuscripts are the primary sources. I saw no URLs for them in search results.

\section{Religious and Community Resources}
I found no religious or community resources.

\section{Digital and Audio Resources}
I found no audio, video or software resources.

\section{Curriculum Development Context}
There are no speakers and no existing programme. The record is a short word list, some sentences and a short text. Teachers could build a small unit on the peoples of the lower Rio Grande and on how extinct languages are known from word lists. They should say plainly how little is known.

\end{document}